% ============================================================
% Section 174: 两端带电均匀棒的转动能级与电场能量修正
% ============================================================
\section{量子力学：两端带电均匀棒的转动能级与电场能量修正}
\label{sec:charged-rod-rotor}

\begin{modelbox}[题目：两端带电均匀棒的转动能级与电场能量修正]
长度为 $d$、质量均匀的棒，可绕其中心在一平面内转动。棒质量为 $M$，两端分别带电荷 $+Q$ 和 $-Q$。
\begin{enumerate}
    \item 写出体系的哈密顿量算符、本征态和本征值；
    \item 若转动平面内存在强度为 $\mathcal{E}$ 的弱电场，用微扰论计算本征值和本征态的变化（波函数精确到一阶，能量精确到二阶）；
    \item 若电场很强，求基态的近似能量。
\end{enumerate}
\end{modelbox}

\begin{tipbox}[考点快速分析]
\begin{itemize}
    \item \textbf{平面转子}：$H_0=-\dfrac{\hbar^2}{2I}\dfrac{d^2}{d\theta^2}$，$I=\dfrac{1}{12}Md^2$，本征态 $e^{im\theta}$
    \item \textbf{微扰}：$H'=-Qd\mathcal{E}\cos\theta$，选择定则 $\Delta m=\pm1$
    \item \textbf{简并子空间}：$m=\pm1$ 需特殊处理（通过 $m=0$ 中间态耦合）
    \item \textbf{强场极限}：$\cos\theta\approx1-\theta^2/2$，化为谐振子
\end{itemize}
\end{tipbox}

\begin{derivationbox}[标准解题过程]
\textbf{1. 无电场时的自由转子}

转动惯量 $I=\dfrac{1}{12}Md^2$，转角 $\theta$。哈密顿量
\[
H_0=-\frac{\hbar^2}{2I}\frac{d^2}{d\theta^2}=-\frac{6\hbar^2}{Md^2}\frac{d^2}{d\theta^2}.
\]

周期性边界条件 $\psi(\theta+2\pi)=\psi(\theta)$，本征函数和本征值：
\begin{equation}
\psi_m^{(0)}(\theta)=\frac{1}{\sqrt{2\pi}}e^{im\theta},\quad E_m^{(0)}=\frac{\hbar^2m^2}{2I}=\frac{6\hbar^2m^2}{Md^2},\quad m=0,\pm1,\pm2,\ldots
\end{equation}

除 $m=0$ 外，能级均为二重简并（$\pm m$）。

\textbf{2. 弱电场中的微扰处理}

设电场沿 $x$ 轴，偶极矩 $p=Qd$（由 $-Q$ 指向 $+Q$），微扰
\[
H'=-Qd\mathcal{E}\cos\theta.
\]

矩阵元 $H'_{m'm}=-\dfrac{Qd\mathcal{E}}{2}(\delta_{m',m+1}+\delta_{m',m-1})$，选择定则 $\Delta m=\pm1$。

\textit{能量一级修正}：对角元为零，$E_m^{(1)}=0$。

\textit{波函数一级修正}（$m\neq\pm1$）：
\[
\psi_m^{(1)}=\frac{Md^3Q\mathcal{E}}{12\hbar^2\sqrt{2\pi}}\Bigl[\frac{e^{i(m+1)\theta}}{2m+1}-\frac{e^{i(m-1)\theta}}{2m-1}\Bigr].
\]

\textit{能量二级修正}（$m\neq\pm1$）：
\begin{equation}
E_m^{(2)}=\frac{MQ^2d^4\mathcal{E}^2}{12(4m^2-1)\hbar^2}.
\end{equation}

\textit{$m=\pm1$ 的简并处理}：

在 $\{|1\rangle,|-1\rangle\}$ 子空间中，二级有效哈密顿矩阵（通过中间态 $|0\rangle$ 耦合）：
\[
H'_{\text{eff}}=\frac{MQ^2d^4\mathcal{E}^2}{72\hbar^2}\begin{pmatrix}2&3\\3&2\end{pmatrix}.
\]

对角化得本征值 $E_{1\pm}^{(2)}=\pm\dfrac{5MQ^2d^4\mathcal{E}^2}{24\hbar^2}$，正确零级波函数为对称/反对称组合 $\dfrac{1}{\sqrt{4\pi}}(e^{i\theta}\pm e^{-i\theta})$。

\textbf{3. 强电场极限}

电场很强时棒沿电场方向排列，$\theta$ 局限于小角度。展开 $\cos\theta\approx1-\theta^2/2$：
\[
H\approx-\frac{\hbar^2}{2I}\frac{d^2}{d\theta^2}+\frac12Qd\mathcal{E}\,\theta^2-Qd\mathcal{E}.
\]

化为谐振子，频率 $\omega=\sqrt{Qd\mathcal{E}/I}$。基态能量和波函数：
\begin{equation}
\boxed{E_0=\frac12\hbar\omega-Qd\mathcal{E}=\frac12\hbar\sqrt{\frac{12Q\mathcal{E}}{Md}}-Qd\mathcal{E},\quad\phi_0(\theta)=\Bigl(\frac{\alpha}{\pi}\Bigr)^{1/4}e^{-\alpha^2\theta^2/2},}
\end{equation}
其中 $\alpha=(Qd\mathcal{E}I/\hbar^2)^{1/4}$。
\end{derivationbox}

\begin{formulabox}[答案汇总]
\begin{center}
\begin{tabular}{ll}
\toprule
\textbf{结论} & \textbf{结果} \\
\midrule
自由转子能级 & $E_m^{(0)}=6\hbar^2m^2/(Md^2)$ \\
弱场二级修正（$m\neq\pm1$） & $E_m^{(2)}=\dfrac{MQ^2d^4\mathcal{E}^2}{12(4m^2-1)\hbar^2}$ \\
$m=\pm1$ 分裂 & $E_{1\pm}^{(2)}=\pm\dfrac{5MQ^2d^4\mathcal{E}^2}{24\hbar^2}$ \\
强场基态能量 & $E_0=\dfrac12\hbar\sqrt{12Q\mathcal{E}/(Md)}-Qd\mathcal{E}$ \\
\bottomrule
\end{tabular}
\end{center}
\end{formulabox}

\begin{insightbox}[物理图像]
\begin{itemize}
    \item 平面转子模型是理解极性分子（如 HCl）在电场中行为的最简单模型。弱场下能量移动为二级 Stark 效应；强场下分子被``锁定''在电场方向，表现为谐振子。
    \item $m=\pm1$ 的简并解除体现了电场对转子取向的偏好：对称组合 $|1\rangle+|-1\rangle\propto\cos\theta$ 对应沿电场方向的摆动（能量下移），反对称组合 $\propto\sin\theta$ 对应垂直于电场的运动（能量上移）。
    \item 该模型也是研究 Josephson 结相位动力学和超导磁通量子比特的理论原型。
\end{itemize}
\end{insightbox}
